\scanpage{102}
\mspage{95}
% archive sheet number: 101
On vérifie aussi immédiatement que
\[
\varinjlim_{\alpha\in\widetilde C}Z_\alpha=\unsure{\underline{\lim}|\varinjlim}X_i=X
\]
donc, $\widetilde C_M$ étant cofinal,
\[
\underline{\lim}_{\alpha\in\widetilde C_M}Z_\alpha=X.
\]
\deleted{Il reste \unsure{bien} \unsure{entendu} que le syst. inductif $Z_\alpha$ sur $\widetilde C$ \unsure{tout entier}, \unsure{restreint|induit} sur $\widetilde C_M$,}

Au total, on est arrivé à représenter $X$ sous la forme
\[
X=\varinjlim_{\alpha\in A}X_\alpha
\]
cette fois les $X_\alpha$ sont \uline{dans $M^{\pi}$} (pas seulement dans $M$), et $A$ est \deleted{\unsure{donc}} \inserted{grand devant $\pi$}. \deleted{adapté à $\pi$ (pas seulement grand devant $\pi$)} Mais on ne sait pas que \deleted{les} toute partie \inserted{filtrante $B$} de cardinal $\leq\pi$ de $A$ admet une borne sup \inserted{$\alpha_B$} dans $A$, et encore moins que $X_{\alpha_B}=\underline{\lim}_{\alpha\in B}X_\alpha$. Donc on ne peut \inserted{encore} appliquer directement b). S'il ne s'agissait que de prouver \inserted{\unsure{la deuxième moitié de|la première moitié de}} a), \unsure{où|ou} on suppose déjà les $X_i\in M^{\pi}$ et $I$ grand devant $I$, on n'a rien gagné du tout par \unsure{nos} \unsure{labeurs}!
% slanted
\margin{Si on \unsure{supprime|suppose} ($\unsure{a_\pi|a^\pi|a}$) cet argument \unsure{nous} prouve \unsure{alors} que $M$ est stable par $\varinjlim$ grandes devant $\pi$.}
